\documentclass{article}
\usepackage[T1]{fontenc}
\usepackage{lmodern}
\usepackage{graphicx}
\usepackage{amsmath,amssymb}
\usepackage[a4paper,margin=2cm]{geometry}
\usepackage[absolute,overlay]{textpos}
\setlength{\TPHorizModule}{1mm}
\setlength{\TPVertModule}{1mm}
\usepackage[indonesian]{babel}
\usepackage{hyperref}
\setlength{\emergencystretch}{2em}
% unit: Halaman 31

\begin{document}

% === Halaman 31 (llm) ===

\section*{B. Cipher Transposisi}

\begin{itemize}
    \item Cipherteks diperoleh dengan mengubah posisi huruf di dalam plainteks.
    \item Dengan kata lain, algoritma ini melakukan \textit{transpose} terhadap rangkaian huruf di dalam plainteks.
    \item Nama lain untuk metode ini adalah \textit{cipher permutasi}, karena \textit{transpose} setiap karakter di dalam teks sama dengan mempermutasikan karakter-karakter tersebut.
    \item Contoh cipher transposisi: \textit{Columnar Transposition Cipher, Rail Fence Transposition Cipher}
\end{itemize}

\end{document}
